\section*{الفصل \ref{ch:b2:curves} --- المنحنيات}

\begin{solution}{exo:b2:curves:1}
$\gamma'(t) = (1 - \cos t,\ \sin t)$، ومنه
\[
\norm{\gamma'(t)}^2 = (1 - \cos t)^2 + \sin^2 t
= 2 - 2\cos t = 4\sin^2\tfrac t2 ,
\]
و$\norm{\gamma'(t)} = 2\sin\frac t2$ (وهو موجب على $[0,
2\pi]$). ومنه
\[
L = \int_0^{2\pi} 2\sin\tfrac t2\,\dd t
= \Bigl[-4\cos\tfrac t2\Bigr]_0^{2\pi} = 8 :
\]
فطول قوس واحد من الدويري هو $8$ (من أجل عجلة نصف قطرها $1$) ---
وهي نتيجة شهيرة لرن، دون أي $\pi$ في الأفق.
\end{solution}

\begin{solution}{exo:b2:curves:2}
مع $x = a\cos t$، $y = b\sin t$: $x' = -a\sin t$، $y' = b\cos
t$، $x'' = -a\cos t$، $y'' = -b\sin t$، ومنه بحسب \cref{prop:b2:curves:kappaformula}
\[
\kappa(t)
= \frac{x'y'' - y'x''}{(x'^2 + y'^2)^{3/2}}
= \frac{ab\sin^2 t + ab\cos^2 t}
       {(a^2\sin^2 t + b^2\cos^2 t)^{3/2}}
= \frac{ab}{(a^2\sin^2 t + b^2\cos^2 t)^{3/2}} .
\]
ويكون المقام أصغريا عندما $\sin t = 0$ (بالقيمة $b^3$، وفي النقطتين
$(\pm a, 0)$) وأعظميا عندما $\cos t = 0$ (بالقيمة $a^3$، وفي النقطتين $(0, \pm b)$)،
لأن $a > b$. ومنه فإن $\kappa$ أعظمي عند طرفي المحور الكبير، $\kappa_{\max} = a/b^2$،
وأصغري عند طرفي المحور الصغير، $\kappa_{\min} = b/a^2$: فالأهليلج أشد انثناء عند
طرفي محوره الطويل.
\end{solution}

\begin{solution}{exo:b2:curves:3}
من أجل التمثيل البياني $\gamma(x) = (x, \cosh x)$: $\norm{\gamma'(x)} =
\sqrt{1 + \sinh^2 x} = \cosh x$، ومنه فإن \omterm{def:b2:curves:length}{طول القوس} من $0$
إلى $x$ هو $\int_0^x \cosh u\,\dd u = \sinh x$. وانحناء تمثيل بياني (\cref{prop:b2:curves:kappaformula}):
\[
\kappa(x) = \frac{f''(x)}{(1 + f'(x)^2)^{3/2}}
= \frac{\cosh x}{\cosh^3 x} = \frac{1}{\cosh^2 x} ,
\]
ومنه $R(x) = \cosh^2 x$، كما أُعلن. ولاحظ التوافق الأنيق $R(x) = 1 + s(x)^2$ مع $s = \sinh x$ \omterm{def:b2:curves:length}{طول
القوس}: فنصف قطر انحناء السلسلية ينمو مع مربع \omterm{def:b2:curves:length}{طول القوس} انطلاقا من
الرأس.
\end{solution}

\begin{solution}{exo:b2:curves:4}
$\gamma'(t) = e^t(\cos t - \sin t,\ \sin t + \cos t)$، ومنه
\[
\langle \gamma(t), \gamma'(t)\rangle = e^{2t}
\bigl(\cos t(\cos t - \sin t) + \sin t(\sin t + \cos t)\bigr)
= e^{2t},
\]
بينما $\norm{\gamma(t)} = e^t$ و$\norm{\gamma'(t)} =
e^t\sqrt 2$. ومنه
\[
\cos\angle\bigl(\gamma, \gamma'\bigr)
= \frac{e^{2t}}{e^t \cdot e^t\sqrt2} = \frac{1}{\sqrt2} :
\]
\omterm{def:b2:curves:arc}{فالمماس} يصنع دائما الزاوية $\pi/4$ مع نصف القطر --- وهي خاصية
\emph{تساوي الزوايا} في الحلزون اللوغاريتمي. \omterm{def:b2:curves:length}{وطول القوس} على $(-\infty, 0]$:
\[
\int_{-\infty}^0 \norm{\gamma'(t)}\,\dd t
= \sqrt2\int_{-\infty}^0 e^t\,\dd t = \sqrt 2 ,
\]
وهو منته مع أن الحلزون يلتف حول المبدأ عددا لانهائيا من المرات.
\omterm{thm:b2:curves:frenet2d}{والانحناء}: مع $x'y'' - y'x''$ المحسوب من $\gamma'' =
e^t(-2\sin t,\ 2\cos t)$،
\[
x'y'' - y'x'' = e^{2t}\bigl(2\cos t(\cos t - \sin t)
+ 2\sin t(\sin t + \cos t)\bigr) = 2e^{2t},
\]
ومنه $\kappa(t) = \dfrac{2e^{2t}}{(e^t\sqrt2)^3} =
\dfrac{1}{e^t\sqrt2}$: \omterm{thm:b2:curves:frenet2d}{فالانحناء} $1/(\sqrt2\,
\norm{\gamma})$، ويتلاشى مع نمو الحلزون.
\end{solution}

\begin{solution}{exo:b2:curves:5}
\emph{القوس الأول:} $\gamma(t) = (t^2,\ t^4 + t^5)$. المشتقات في $0$: $\gamma'' = (2, 0) \neq 0$، ومنه $p = 2$.
ثم $\gamma^{(3)}(0)
= (0, 0)$، $\gamma^{(4)}(0) = (0, 24)$، وهي غير خطية مع $(2,
0)$: $q = 4$. وكلاهما زوجي:
\emph{نقطة رجوع من النوع الثاني} --- فالفرعان يغادران في الاتجاه
$+u = (1,0)$ ويبقيان في الجهة نفسها من \omterm{def:b2:curves:arc}{المماس}. (وفعلا $y = x^2 \pm x^{5/2}$ على الفرعين:
بالإشارة نفسها من أجل $x$ صغيرة.)

\emph{القوس الثاني:} $\gamma(t) = (t^3, t^4)$. $\gamma'(0) =
\gamma''(0) = 0$، $\gamma^{(3)}(0) = (6, 0)$: $p = 3$، وهو فردي. ثم
$\gamma^{(4)}(0) = (0, 24)$: $q = 4$، وهو زوجي. فردي--زوجي: \emph{نقطة عادية} --- فرغم
انعدام السرعة، يعبر المسار $y = x^{4/3}$ المبدأ بسلاسة، ويبقى فوق مماسه
$y = 0$.
\end{solution}

\begin{solution}{exo:b2:curves:6}
نفاضل $c(s) = \gamma(s) + \dfrac{1}{\kappa(s)}N(s)$ باستعمال صيغ فرينيه في المستوي (\cref{thm:b2:curves:frenet2d}):
\[
c'(s) = T(s) - \frac{\kappa'(s)}{\kappa(s)^2}N(s)
+ \frac{1}{\kappa(s)}\,\bigl(-\kappa(s)T(s)\bigr)
= -\frac{\kappa'(s)}{\kappa(s)^2}\,N(s) ,
\]
مع تلاشي الحدود المماسية بالضبط. ومنه فإن سرعة المطوّرة يحملها
المتجهة $N(s)$، وهي توجه ناظم $\gamma$ في $\gamma(s)$ --- والنقطة $c(s)$ تقع
على ذلك الناظم نفسه: أي إن المطوّرة هي \emph{مغلِّف النواظم}. (وحيث
يكون $\kappa' = 0$ تكون للمطوّرة نقطة شاذة؛ وهذا ما ينتج نقط رجوع مطوّرة
الأهليلج.)
\end{solution}

\begin{solution}{exo:b2:curves:7}
لنضع $G(s) = F(s)F(s)^{\mathsf T}$. عندئذ، باستعمال $F' = FA$ و$(F^{\mathsf T})' = (F')^{\mathsf T} = A^{\mathsf T}F^{\mathsf T}$،
\[
G' = F'F^{\mathsf T} + F(F^{\mathsf T})'
= FAF^{\mathsf T} + FA^{\mathsf T}F^{\mathsf T}
= F(A + A^{\mathsf T})F^{\mathsf T} = 0
\]
بحكم تخالف التناظر. ومنه فإن $G$ ثابتة على الفترة، وتساوي $G(s_0) = F(s_0)F(s_0)^{\mathsf T} = I$:
أي إن $F(s)$ متعامدة من أجل كل $s$. (وبديلا عن ذلك، دون حساب
$G'$ وإثبات انعدامه: تحل كل من $G$ والمصفوفة الثابتة $I$
الجملة الخطية $Y' = YA +
A^{\mathsf T}Y$ بالقيمة الابتدائية نفسها، وتفرض وحدانية
كوشي--ليبشيتز للجمل الخطية، \cref{ch:b2:diffeq}، أن $G
\equiv I$.)

\emph{الصلة:} لجملة فرينيه $(T, N, B)' = (T, N, B)\,A(s)$ مصفوفة معاملات متخالفة التناظر
\[
A = \begin{pmatrix} 0 & -\kappa & 0\\ \kappa & 0 & -\tau\\
0 & \tau & 0\end{pmatrix}
\]
(وأعمدتها تعبّر عن $T', N', B'$). ويبين الحساب أعلاه أن معلم حل ينطلق
متعامدا \emph{يبقى} متعامدا --- وهي الخطوة المفتاحية في المبرهنة
الأساسية التي تعيد بناء منحن انطلاقا من $(\kappa, \tau)$.
\end{solution}

\begin{solution}{exo:b2:curves:8}
بحسب \cref{thm:b2:curves:fundamental} (شق الوحدانية)، توجد دالة زاوية $\varphi$ من الصنف $\mathcal{C}^1$
تحقق $T(s) =
(\cos\varphi(s), \sin\varphi(s))$ و$\varphi' = \kappa$. ومنه
\[
\int_0^L \kappa(s)\,\dd s = \varphi(L) - \varphi(0) .
\]
وبما أن القوس مغلق ودوره $L$، فإن $T(L) = T(0)$: $(\cos\varphi(L), \sin\varphi(L)) = (\cos\varphi(0),
\sin\varphi(0))$، ومنه $\varphi(L) - \varphi(0) \in 2\pi\Z$.
ومنه فإن \omterm{thm:b2:curves:frenet2d}{الانحناء} الكلي لمنحن مغلق مضاعف صحيح للمقدار $2\pi$ دائما
--- والعدد الصحيح هو \emph{عدد لف} \omterm{def:b2:curves:arc}{المماس} (أي عدد الدورات الكاملة
التي يصنعها $T$). ومن أجل الدائرة ذات نصف القطر $R$: $\kappa = 1/R$
و$L = 2\pi R$، \omterm{thm:b2:curves:frenet2d}{فالانحناء} الكلي $2\pi$ وعدد اللف $1$؛ وتنص مبرهنة
المماسات الدوارة على أن هذه القيمة تصلح لكل منحن مغلق بسيط.
\end{solution}

\begin{solution}{exo:b2:curves:9}
بما أن $\tau \equiv 0$، فإن المنحنى يقع في مستو (\cref{prop:b2:curves:torsion})؛ فلنعمل في ذلك
المستوي. ولننظر في المركز المرشح
\[
c(s) = \gamma(s) + \frac{1}{\kappa}N(s)
\qquad (\kappa \text{ ثابت}).
\]
وبالمفاضلة بصيغ فرينيه (فهنا $N' = -\kappa T + \tau B
= -\kappa T$):
\[
c'(s) = T + \frac1\kappa(-\kappa T) = 0 ,
\]
ومنه فإن $c$ نقطة ثابتة $\Omega$. ثم $\norm{\gamma(s) -
\Omega} = \norm{-\frac1\kappa N(s)} = \frac1\kappa$ من أجل كل $s$:
فالمنحنى يقع على الدائرة ذات المركز $\Omega$ ونصف القطر $1/\kappa$ (في
مستويه)، وبما أنه قوس غير ثابت منها، فهو قوس من تلك الدائرة.
\end{solution}

\begin{solution}{exo:b2:curves:10}
من أجل التمثيل البياني $\gamma(x) = (x, x^2/2)$، $\norm{\gamma'(x)} =
\sqrt{1 + x^2}$، ومنه $L = \int_0^a\sqrt{1+x^2}\,\dd x$.
وبالتعويض $x = \sinh u$ ($\dd x = \cosh u\,\dd u$، و$u$ من $0$ إلى $u_a = \ln(a + \sqrt{1+a^2})$):
\[
L = \int_0^{u_a}\cosh^2 u\,\dd u
= \frac12\bigl[u + \sinh u\cosh u\bigr]_0^{u_a}
= \frac12\Bigl(\ln\bigl(a + \sqrt{1+a^2}\bigr) +
a\sqrt{1+a^2}\Bigr),
\]
باستعمال $\cosh^2 u = \frac{1 + \cosh 2u}2$ و$\sinh u_a = a$، $\cosh u_a = \sqrt{1 + a^2}$.
\end{solution}

\begin{solution}{exo:b2:curves:11}
نوسّم \omterm{def:b2:curves:length}{بطول القوس} (\cref{thm:b2:curves:arclength}) ونضع $\lambda(s) =
\langle P - \gamma(s), T(s)\rangle$، وهي دالة من الصنف $\mathcal C^1$؛
وبما أن $P$ تقع على \omterm{def:b2:curves:arc}{المماس} في $\gamma(s)$، فإن المتجهة $P - \gamma(s)$ خطية مع
$T(s)$، ومنه $P =
\gamma(s) + \lambda(s)T(s)$. وبالمفاضلة،
\[
0 = T(s) + \lambda'(s)T(s) + \lambda(s)T'(s)
= \bigl(1 + \lambda'(s)\bigr)T(s) + \lambda(s)T'(s),
\]
و$T'(s) \perp T(s)$ (بمفاضلة $\norm T^2 = 1$)، ومنه تنعدم المركبتان معا: $\lambda' = -1$ و$\lambda T' = 0$.
ثم إن $\lambda(s) = c - s$ ينعدم مرة واحدة على الأكثر، ومنه $T' = 0$ على مجموعة
كثيفة، ومنه في كل نقطة بالاتصال: فالمتجهة $T$ واحدية ثابتة
و$\gamma(s) = \gamma(s_0) + (s -
s_0)T$: أي مستقيم (يمر بالنقطة $P$، كما يجب).
\end{solution}

\begin{solution}{exo:b2:curves:12}
(1) مصفوفة فرينيه
\[
A(s) = \begin{pmatrix} 0 & -\kappa & 0\\ \kappa & 0 & -\tau\\
0 & \tau & 0\end{pmatrix}
\]
معاملاتها متصلة، ومنه فإن الجملة الخطية $F' = FA(s)$، $F(s_0) = F_0$ (بمصفوفة
متعامدة مباشرة) لها حل وحيد على $J$ كله (\cref{thm:b2:diffeq:linear}). وبحسب \cref{exo:b2:curves:7}،
تكون $F(s)$ متعامدة من أجل كل $s$؛ والمقدار $\det F$ \omterm{def:b2:metric:continuity}{متصل} بقيم
في $\{\pm1\}$ ويساوي $1$ في $s_0$، ومنه فإن $F(s)$ مباشرة من أجل
كل $s$.

(2) نقرأ الأسطر $T, N, B$ للمصفوفة $F$ (بحيث $T' = \kappa
N$، $N' = -\kappa T + \tau B$،
$B' = -\tau N$) ونضع $\gamma(s) = \gamma_0 + \int_{s_0}^s T(u)\,\dd u$. عندئذ يكون $\gamma' = T$ متجهة واحدية: أي السرعة
الواحدية؛ و$T' = \kappa N$ مع $\kappa > 0$ و$N$ واحدي ومتعامد مع $T$، ومنه
فإن $\gamma$ ثنائي النظامية وانحناؤه $\norm{T'} = \kappa$ وناظمه الرئيسي $N$؛
والناظم الثنائي هو $T \wedge N = B$ (فالمعلم متعامد مباشر)، وتحدد $B' = -\tau N$
\omterm{thm:b2:curves:frenet3d}{الالتواء} بالمقدار $\tau$.

(3) ليكن $\gamma_1, \gamma_2$ منحنيين ذوي سرعة واحدية وثنائيي النظامية لهما
$(\kappa, \tau)$ نفسه. يوجد تقايس مباشر وحيد $\Phi = \rho + w$ ($\rho \in SO(3)$) يرسل $\gamma_1(s_0)$ على
$\gamma_2(s_0)$ ويرسل معلم فرينيه للمنحنى $\gamma_1$ في $s_0$ على معلم $\gamma_2$
في $s_0$. والمنحنى $\Phi\circ\gamma_1$ ذو سرعة واحدية وله الثوابت نفسها
(فمعلمه هو $\rho$ مطبقا على معلم $\gamma_1$، والتطبيق $\rho$ يحافظ على
الجداءات الشعاعية لأنه مباشر). والآن يحل معلما $\Phi\circ\gamma_1$ و$\gamma_2$ معا
الجملة $F' =
FA(s)$ بالقيمة الابتدائية نفسها، ومنه فهما متطابقان
بالوحدانية؛ وبوجه خاص يتطابق المماسان، وبالمكاملة انطلاقا من النقطة
المشتركة $s_0$: $\Phi\circ\gamma_1 = \gamma_2$.
\end{solution}

\begin{solution}{pb:b2:curves:1}
\textbf{1.} الجملة خطية بدلالة $(x, y)$ ومحددها $\Delta(t) = a b' - a'b \neq 0$: وتعطي قاعدة
كرامر الحل الوحيد
\[
x = \frac{c b' - c' b}{\Delta}, \qquad
y = \frac{a c' - a' c}{\Delta}.
\]

\textbf{2.} بما أن $a(t)x(t) + b(t)y(t) = c(t)$ تطابقا، فإن المفاضلة تعطي $a'x + b'y + ax' + by' = c'$؛
وتقتل المعادلة المميزة الثانية المقدار $a'x + b'y - c'$، ومنه $a(t)x'(t) + b(t)y'(t) = 0$: أي إن
$E'(t)$ متعامد مع $(a,
b)$، ومنه فهو موازٍ للمتجهة $(-b, a)$، وهو اتجاه
$D_t$. وبما أن $E(t) \in D_t$ (بالمعادلة الأولى) و$E'(t) \neq 0$، فإن المستقيم
$D_t$ هو بالضبط \omterm{def:b2:curves:arc}{مماس} المنحنى $E$ في $E(t)$.

\textbf{3.} هنا $(a, b, c) = (t, -1, t^2/2)$، ومنه $\Delta =
t\cdot0 - 1\cdot(-1) = 1$ و
\[
x = \frac{c b' - c' b}{\Delta} = \tfrac{t^2}2\cdot 0 +
t = t, \qquad
y = \frac{a c' - a' c}{\Delta} = t\cdot t - \tfrac{t^2}2 =
\tfrac{t^2}2 :
\]
فمغلِّف مماسات القطع المكافئ هو القطع المكافئ، كما يجب.

\textbf{4.} الناظم هو $\langle M, T(s)\rangle =
\langle\gamma(s), T(s)\rangle$: بالمعاملات $a = T_1$ و$b =
T_2$ و$c = \langle\gamma, T\rangle$.
وبالمفاضلة بصيغ فرينيه ($T' = \kappa N$): $a' = \kappa N_1$، $b' = \kappa
N_2$، و$c' = \langle T, T\rangle + \langle\gamma, \kappa
N\rangle = 1 + \kappa\langle\gamma, N\rangle$. وتُكتب
المعادلة المميزة الثانية $\kappa\langle M, N\rangle = 1 +
\kappa\langle\gamma, N\rangle$ على الصورة $\kappa\langle M -
\gamma, N\rangle = 1$. وتقول الأولى
$M - \gamma \perp T$، ومنه $M - \gamma = \mu N$ مع $\mu = 1/\kappa$: فالنقطة المميزة هي $\gamma + \frac1\kappa N$، أي
\omterm{def:b2:curves:curvature}{مركز الانحناء}، ومغلِّف النواظم هو مطوّرة \cref{exo:b2:curves:6}. وأخيرا $\Delta = T_1\kappa N_2 -
\kappa N_1 T_2 = \kappa\det(T, N) = \kappa \neq 0$.

\textbf{5.} الحزمة: للجملة $x\cos\theta + y\sin\theta =
0$، $-x\sin\theta + y\cos\theta = 0$ \omterm{def:b2:linalg:det}{محدد} $1$ والحل $(0,0)$
من أجل كل $\theta$: $E \equiv (0,0)$، $E'
\equiv 0$، ولا وجود لمنحن --- بل مجرد
النقطة المشتركة بين جميع المستقيمات. والعائلة المتوازية: يعطي $a' = b' = 0$
المقدار $\Delta \equiv 0$، والمعادلة الثانية $0 = c'(t)$، وهي تسقط بمجرد أن تتحرك
العائلة فعلا: فلا نقطة مميزة، وفعلا لا تلامس عائلة مستقيمات متوازية
أي منحن على طول جميع عناصرها.

\textbf{6.} المستقيم المار بالنقطتين $(\cos t, 0)$ و$(0, \sin t)$ هو $\frac x{\cos t} + \frac y{\sin t} = 1$، أي
$x\sin t +
y\cos t = \sin t\cos t$. ومع $(a, b, c) = (\sin t, \cos t,
\sin t\cos t)$: $a' = \cos t$، $b' = -\sin t$، $c' = \cos
2t$، $\Delta = -\sin^2 t - \cos^2 t = -1$. وبقاعدة
كرامر:
\begin{align*}
x &= \frac{cb' - c'b}{-1} = \sin^2 t\cos t + \cos 2t\cos t
= \cos t\,(\sin^2 t + \cos^2 t - \sin^2 t) = \cos^3 t,\\
y &= \frac{ac' - a'c}{-1} = \sin t\cos^2 t - \sin t\cos 2t
= \sin t\,(\cos^2 t - \cos^2 t + \sin^2 t) = \sin^3 t .
\end{align*}
ثم $(\cos^3t)^{2/3} + (\sin^3t)^{2/3} = 1$: وهي \omterm{pb:b2:curves:1}{النجمية}.

\textbf{7.} $E'(t) = 3(-\cos^2 t\sin t,\ \sin^2 t\cos t)$ ينعدم عند $t = 0$. وهناك يعطي $E''(0) = (-3, 0) \neq 0$ المقدار
$p = 2$؛ ومركبة $E$ على المحور $x$ زوجية بدلالة $t$، ومنه
$E'''(0) = (0, 6)$، وهي غير خطية معه: $q = 3$. زوجي--فردي: أي نقطة رجوع من النوع
الأول في $(1, 0)$ (\cref{prop:b2:curves:local})، \omterm{def:b2:curves:arc}{بمماس} موازٍ للمحور $x$. وتنقل تناظرات
\omterm{pb:b2:curves:1}{النجمية} $x \mapsto -x$ و$y \mapsto -y$ و$(x,
y) \mapsto (y, x)$ نقطة الرجوع إلى $(-1,
0)$ و$(0, \pm1)$.

\textbf{8.} $E'(t) = 3\sin t\cos t\,(-\cos t, \sin t)$، ومنه $\norm{E'(t)} = 3\abs{\sin t\cos t} = \tfrac32\abs{\sin 2t}$.
وفي ربع واحد: $\int_0^{\pi/2}\tfrac32\sin 2t\,\dd t =
\tfrac32$، وبالتناظر يكون الطول الكلي $4 \cdot
\tfrac32 = 6$.

\textbf{9.} $E(t) - P_t = (\cos^3 t - \cos t,\ \sin^3 t) =
\sin^2 t\,(-\cos t,\ \sin t) = \sin^2 t\,(Q_t - P_t)$. ومنه فإن نقطة التماس هي مركز ثقل $(P_t, \cos^2 t)$ و$(Q_t, \sin^2 t)$:
فعندما يجري $t$ من $0$ إلى $\pi/2$ تنزلق من طرف السلم الأرضي
إلى طرفه الجداري.

\textbf{10.} في الربع الأول، مساحة المنطقة الواقعة تحت \omterm{pb:b2:curves:1}{النجمية} هي
$\int_0^1 y\,\dd x$ مع $x = \cos^3 t$ متناقصة من $1$ إلى $0$ عندما يجري $t$
من $0$ إلى $\pi/2$:
\[
\int_0^1 y\,\dd x
= \int_{\pi/2}^{0}\sin^3 t\,(-3\cos^2 t\sin t)\,\dd t
= 3\int_0^{\pi/2}\sin^4 t\cos^2 t\,\dd t .
\]
وبالتخطيط: $\sin^4 t\cos^2 t = (\sin t\cos t)^2\sin^2 t =
\tfrac18\bigl(\sin^2 2t - \sin^2 2t\cos 2t\bigr)$، و$\int_0^{\pi/2}\sin^2 2t\,\dd t = \tfrac\pi4$ بينما $\int_0^{\pi/2}\sin^2 2t\cos 2t\,\dd t =
\bigl[\tfrac{\sin^3 2t}6\bigr]_0^{\pi/2} = 0$. ومنه فإن التكامل $\tfrac\pi{32}$،
ومساحة الربع $\tfrac{3\pi}{32}$، والمساحة المحصورة $4 \cdot
\tfrac{3\pi}{32} = \tfrac{3\pi}8$.

\textbf{11.} \omterm{def:b2:curves:arc}{المماس} في $\gamma(t) = (t, t^2/2)$ موجه بالمتجهة $(1, t)$، ومنه فإن الناظم
هو $\{(x, y) : (x -
t) + t\,(y - t^2/2) = 0\}$، أي $x + t\,y = t +
\tfrac{t^3}2$.

\textbf{12.} $(a, b, c) = (1, t, t + t^3/2)$: $a' = 0$، $b' =
1$، $c' = 1 + \tfrac32 t^2$، $\Delta = 1$. ثم $y = ac' -
a'c = 1 + \tfrac32t^2$ و
\[
x = cb' - c'b = t + \tfrac{t^3}2 - t\Bigl(1 +
\tfrac32t^2\Bigr) = -t^3 .
\]
وبحذف $t$: $t^2 = \tfrac23(y - 1)$ و$x^2 = t^6 =
\tfrac8{27}(y - 1)^3$: أي قطع مكافئ نصف تكعيبي رأسه $(0, 1)$.

\textbf{13.} $T = (1, t)/\sqrt{1+t^2}$، $N = (-t,
1)/\sqrt{1+t^2}$، و$\kappa = (1+t^2)^{-3/2}$، ومنه
\[
\gamma + \frac1\kappa N = (t, \tfrac{t^2}2) +
(1+t^2)\,(-t, 1) = \Bigl(-t^3,\ 1 + \tfrac32t^2\Bigr),
\]
أي المنحنى نفسه: فمغلِّف النواظم هو محل مراكز \omterm{thm:b2:curves:frenet2d}{الانحناء}، كما وعد
السؤال 4.

\textbf{14.} $E'(t) = (-3t^2, 3t)$ ينعدم عند $t = 0$؛ ويعطي $E''(0) = (0, 3) \neq 0$ المقدار $p = 2$،
ويعطي $E'''(0) = (-6,
0)$ المقدار $q = 3$: أي نقطة رجوع من النوع الأول في $(0, 1)$.
والرأس هو حيث يكون $\kappa = (1+t^2)^{-3/2}$ أعظميا، ومنه $\kappa'(0) = 0$، وتنعدم سرعة المطوّرة
$-\frac{\kappa'}{\kappa^2}N$ هناك بالضبط: فنقط رجوع المطوّرة تقع عند القيم القصوى
للانحناء.

\textbf{15.} يمر الناظم عند الوسيط $t$ بالنقطة $(x_0, y_0)$ إذا وفقط
إذا كان $x_0 + t\,y_0 = t + \tfrac{t^3}2$، أي
\[
\frac{t^3}2 + (1 - y_0)\,t - x_0 = 0 ,
\]
وهي معادلة تكعيبية بدلالة $t$: فلها جذر حقيقي واحد أو ثلاثة
(دون حساب التعدد، من أجل النقط العامة). وعلى المحور $x_0 = 0$ تتفكك
على الصورة $t\bigl(\tfrac{t^2}2 + 1 - y_0\bigr) = 0$: أي الجذر $t = 0$ (فالمحور هو الناظم عند الرأس)،
مضافا إليه $t = \pm\sqrt{2(y_0 - 1)}$ عندما $y_0 > 1$. ومنه: ناظم واحد من أجل $y_0 < 1$،
وثلاثة من أجل $y_0 > 1$، وعند $y_0 = 1$ يدل الجذر الثلاثي على نقطة رجوع
المطوّرة. وبوجه عام، يعني جذر مضاعف للمعادلة التكعيبية أن النقطة
تحقق معادلة المستقيم ومشتقتها بالنسبة إلى $t$ معا --- أي إنها
تقع \emph{على المغلِّف}: فالمطوّرة هي بالضبط الحد الفاصل بين منطقة
الناظم الواحد ومنطقة النواظم الثلاثة.

\textbf{16.} يعكس الانعكاس في المرآة المركبة الناظمية للاتجاه ويبقي
المركبة المماسية: فبكتابة $u = \langle u, n\rangle n + u_{\mathrm{tan}}$، يكون الاتجاه المنعكس $u_{\mathrm{tan}} - \langle u,
n\rangle n = u - 2\langle u, n\rangle n$. وهنا
$\langle u,
n\rangle = \cos\theta$، ومنه
\[
v = (1, 0) - 2\cos\theta\,(\cos\theta, \sin\theta)
= (1 - 2\cos^2\theta,\ -2\sin\theta\cos\theta)
= -(\cos2\theta,\ \sin2\theta).
\]

\textbf{17.} يمر الشعاع المنعكس بالنقطة $P_\theta =
(\cos\theta, \sin\theta)$ باتجاه $(\cos2\theta,
\sin2\theta)$؛
والمتجهة الناظمية له هي $(-\sin2\theta,
\cos2\theta)$، ومنه فإن المستقيم هو
\[
-\sin2\theta\,(x - \cos\theta) + \cos2\theta\,(y -
\sin\theta) = 0,
\]
والثابت هو $-\sin2\theta\cos\theta +
\cos2\theta\sin\theta = -\sin\theta$: وبالضرب في $-1$، $x\sin2\theta - y\cos2\theta = \sin\theta$.

\textbf{18.} $(a, b, c) = (\sin2\theta, -\cos2\theta,
\sin\theta)$: $a' = 2\cos2\theta$، $b' = 2\sin2\theta$، $c' =
\cos\theta$، $\Delta = 2\sin^22\theta + 2\cos^22\theta = 2$.
وبقاعدة كرامر، ثم بصيغ تحويل الجداء إلى مجموع:
\begin{align*}
x &= \frac{2\sin\theta\sin2\theta +
\cos\theta\cos2\theta}{2}
= \frac{(\cos\theta - \cos3\theta) + \frac12(\cos\theta +
\cos3\theta)}{2} = \frac{3\cos\theta - \cos3\theta}4,\\
y &= \frac{\sin2\theta\cos\theta - 2\cos2\theta\sin\theta}2
= \frac{\frac12(\sin3\theta + \sin\theta) - (\sin3\theta -
\sin\theta)}2 = \frac{3\sin\theta - \sin3\theta}4 :
\end{align*}
وهي \omterm{pb:b2:curves:1}{الكلوية}، منحن مغلق بنقطتي رجوع.

\textbf{19.} بالمفاضلة والتفكيك مع $\sin3\theta
- \sin\theta = 2\cos2\theta\sin\theta$ و$\cos\theta -
\cos3\theta = 2\sin2\theta\sin\theta$:
\[
E'(\theta) = \tfrac34\bigl(\sin3\theta - \sin\theta,\
\cos\theta - \cos3\theta\bigr)
= \tfrac32\sin\theta\,(\cos2\theta, \sin2\theta),
\]
وهو موازٍ للاتجاه المنعكس في السؤال 16: فكل شعاع منعكس \omterm{def:b2:curves:arc}{مماس} للمنحنى
الحارق، كما تقتضي خاصية المغلِّف. وينعدم $E' = 0$ عند $\sin\theta = 0$ بالضبط:
$E(0) = (\tfrac12, 0)$ و$E(\pi) = (-\tfrac12, 0)$، وهما نقطتا الرجوع. وبوضع $y = 0$ في معادلة المستقيم:
$x\sin2\theta =
\sin\theta$، ومنه $x = \frac1{2\cos\theta} \to \frac12$ عندما $\theta \to 0$: فالأشعة المحاذية للمحور تتجمع
على المسافة $R/2$ من المركز --- وهو البعد البؤري للمرآة الكروية.

\textbf{20.} $\norm{E'(\theta)} = \tfrac32\abs{\sin\theta}$، ومنه فإن الطول $\tfrac32\int_0^{2\pi}\abs{\sin\theta}\,
\dd\theta = \tfrac32\cdot4 = 6$. وبجوار $\theta = 0$، مع
$\cos k\theta = 1 - \tfrac{k^2\theta^2}2 + O(\theta^4)$ و$\sin k\theta = k\theta - \tfrac{k^3\theta^3}6 +
O(\theta^5)$:
\[
x - \tfrac12 = \frac{3\theta^2 + O(\theta^4)}{4}
= \tfrac34\theta^2 + O(\theta^4),
\qquad
y = \frac{4\theta^3 + O(\theta^5)}{4}
= \theta^3 + O(\theta^5) :
\]
$p = 2$، $q = 3$: أي نقطة رجوع من النوع الأول متجهة على طول المحور
--- وهي النقطة المضيئة في المنحنى الحارق لفنجان القهوة.

\textbf{21.} نوسّم النقط المضاءة بالمقدار $\Phi(\theta, r) = P_\theta + r\,v_\theta$ (وهو الموضع على
طول كل شعاع منعكس). \omterm{def:b2:linalg:det}{والمحدد} الجاكوبي $\det(\partial_\theta\Phi, \partial_r\Phi) =
\det(P_\theta' + rv_\theta', v_\theta)$ أفيني بدلالة $r$
وينعدم من أجل $r = r_*(\theta)$ واحد بالضبط --- والنقطة $\Phi(\theta, r_*(\theta))$ هي النقطة
المميزة، لأن اتجاه الشعاع وتغير العائلة يصيران هناك مرتبطين. وخارج
المغلِّف يكون التطبيق تقابلا تفاضليا محليا، والنقطة الواقعة داخل
المنحنى الحارق مباشرة يصيبها شعاعان متجاوران (فهناك حلان $\theta$)،
أما النقطة الخارجة فلا يصيبها أي شعاع من ذلك الجزء من العائلة؛ وعلى
المنحنى الحارق يندمج الشعاعان. وشدة الضوء عكسية التناسب مع القيمة
المطلقة \omterm{def:b2:linalg:det}{للمحدد} الجاكوبي، ومنه فهي تنفجر على طول المغلِّف: فالمنحنى
الحارق هو المنحنى المضيء، وأشد ما يكون إضاءة عند نقطة الرجوع، حيث
يبلغ الانحلال أقصاه.

\textbf{22.} $x' = 1 - \cos t$، $y' = \sin t$، $x'' = \sin
t$، $y'' = \cos t$، ومنه $x'y'' - y'x'' = \cos t - 1$ و$\norm{\gamma'}^2 = 2(1 - \cos t) = 4\sin^2\tfrac t2$؛
وبحسب \cref{prop:b2:curves:kappaformula}،
\[
\kappa(t) = \frac{-(1 - \cos t)}{8\sin^3\tfrac t2}
= \frac{-2\sin^2\tfrac t2}{8\sin^3\tfrac t2}
= -\frac1{4\sin\tfrac t2}
\qquad (0 < t < 2\pi).
\]
ومع $T = (\sin\tfrac t2, \cos\tfrac t2)$ (بقسمة $\gamma'$ على $2\sin\tfrac t2$) و$N = (-\cos\tfrac t2, \sin\tfrac
t2)$:
\[
\gamma + \frac1\kappa N
= \gamma - 4\sin\tfrac t2\,\Bigl(-\cos\tfrac t2,\
\sin\tfrac t2\Bigr)
= \bigl(t - \sin t + 2\sin t,\ 1 - \cos t - 2(1 - \cos
t)\bigr),
\]
أي $c(t) = (t + \sin t,\ \cos t - 1)$. وبالتعويض $t = u
+ \pi$:
\[
c = \bigl(u + \pi - \sin u,\ -\cos u - 1\bigr)
= \bigl((u - \sin u) + \pi,\ (1 - \cos u) - 2\bigr) :
\]
أي الدويري $\gamma(u)$ منسحبا بالمتجهة $(\pi, -2)$. فمطوّرة الدويري دويري
مطابق، معلق مستوى واحدا تحته.

\textbf{23.} $R(t) = 1/\abs{\kappa(t)} = 4\sin\tfrac t2$، ومنه $R(\pi) = 4$: أي نصف \omterm{def:b2:curves:length}{طول القوس} $8$ المحسوب
في \cref{exo:b2:curves:1}. وتنعدم سرعة المطوّرة $c'(t) = (1 +
\cos t, -\sin t)$ عند $t = \pi$: فنقطة الرجوع هي
$c(\pi)
= (\pi, -2)$، أسفل القمة $\gamma(\pi) = (\pi,
2)$ مباشرة، على المسافة $4 = R(\pi)$، وهي بالضبط
طول نصف قطر التقبيل هناك.

\textbf{24.} انطلاقا من \cref{exo:b2:curves:6}، $c'(s) =
-\frac{\kappa'(s)}{\kappa(s)^2}N(s) = R'(s)\,N(s)$، ومنه $\norm{c'(s)} = \abs{R'(s)}$، ومن أجل
$R$ رتيبة،
\[
\int_{s_0}^{s_1}\norm{c'(s)}\,\dd s =
\Bigl|\int_{s_0}^{s_1}R'(s)\,\dd s\Bigr| = \abs{R(s_1) -
R(s_0)} .
\]
ولنقل إن $R$ متناقصة. فخيط موضوع على طول المطوّرة إلى ما بعد
$c(s_0)$ ومطوّل بالقطعة الواصلة من $c(s_0)$ إلى $\gamma(s_0)$ (وهي مماسة
للمطوّرة، بحسب السؤال 4) يكون له، إذا قُشر إلى غاية $c(s)$ وشُدّ،
جزء مستقيم طوله $R(s_0) - \bigl(R(s_0) - R(s)\bigr) = R(s)$ ومتجه من $c(s)$ على طول الناظم --- فيحط
بالضبط على $\gamma(s)$: أي إن الطرف الحر يرسم المنحنى الأصلي (وهو
``المنحنى المنشور''). وقد علّق هويغنز نوّاسا بين خدّين دويريين:
فينلف الخيط على المطوّرة، ومنه يرسم الثقل دويريا --- وهو المنحنى
المتساوي الزمن، الذي لا يتعلق دور تذبذبه بالسعة.

\textbf{25.} مماسات القطع المكافئ: $\Delta = 1$ و$E' =
(1, t) \neq 0$ في كل نقطة ---
فالمغلِّف أملس (وهو القطع المكافئ نفسه). السلم المنزلق: $\Delta = -1$، لكن
$E' =
\tfrac32\sin 2t\,(-\cos t, \sin t)$ ينعدم عند طرفي الربع --- وهي نقط الرجوع الأربع للنجمية.
نواظم القطع المكافئ ونواظم الدويري: $\Delta = \kappa \neq 0$، وينعدم $E' = R'N$ حيث يكون
\omterm{thm:b2:curves:frenet2d}{الانحناء} أقصويا بالضبط --- فنقط رجوع المطوّرتين عند $(0,1)$ و$(\pi, -2)$.
المنحنى الحارق: $\Delta = 2$، وينعدم $E' = \tfrac32\sin\theta\,(\cos2\theta,
\sin2\theta)$ عند $\theta = 0, \pi$ --- وهما نقطتا
رجوع \omterm{pb:b2:curves:1}{الكلوية}، أي بؤرتا المرآة. والقيم القصوى للانحناء \emph{يجب} أن
تنتج نقط رجوع على مغلِّف النواظم، لأن سرعة المطوّرة هي $R'N$: ولهذا
فإن لمطوّرة الأهليلج أربع نقط رجوع (أربعة رؤوس)، والعائلتان
المنحلتان (الحزمة والمتوازيات) هما الحالتان اللتان تُخرج فيهما
الآلة نقطة أو لا شيء البتة.
\end{solution}
